%% kapa: exam=1 round=26 year=2015 period=1 form=A subject=civil_law subject_ko=민법 seq=28 number=28
%% source: Q-Net 감정평가사 기출문제 (cst00309) artlSeq=5138001 file=extracted/1008/5138001/제26회_감정평가사 1차_기출/제26회_감정평가사 1차_기출/1교시_감평_A형.pdf
%% conversion: kordoc 4.24.51 -> tools/gen1.py (auto); review=proofread-sonnet; flags=none
\begin{question}{28}
「민법」상 공동소유에 관한 설명으로 옳지 않은 것은? (특약은 고려하지 않음)
\begin{choices}
  \item 공유물의 임대는 공유자의 과반수로 결정한다.
  \item 공유자는 내부적 관계에서 지분의 비율로 공유물의 관리비용 기타 의무를 부담한다.
  \item 조합재산이 아닌 합유물을 처분하기 위해서는 합유자 전원의 동의가 있어야 한다.
  \item 합유물의 지분을 처분하기 위해서는 합유자 전원의 동의가 있어야 한다.
  \item 총유물의 처분은 물론 관리도 사원총회의 결의에 의해야 한다.
\end{choices}
\end{question}
\answer{1}
